% Dense in this paper != dense as in rationals are dense in reals.
% When we say low dimensional, we mean relative to the size of the vocabulary. So a 300-dimensional embedding for a vocabulary of 10,000 words is low dimensional.

\thispagestyle{plain}
\begin{center}
  \large{MERKINTÄTAVOISTA}
\end{center}
Tässä opinnäytetyössä käytetään seuraavia merkintöjä ja termejä.

\begin{align*}
  &\mathbb{R} &&\text{reaalilukujen joukko} \\
  &\mathbb{N} &&\text{luonnollisten lukujen joukko } \{1, 2, 3, \dots\} \\
  &A^{\top} &&\text{matriisin tai vektorin } A \text{ transpoosi} \\
  &\odot &&\text{Hadamard-tulo (alkioittainen kertolasku)} \\
  &[A \,\|\, B] &&\text{vektorien tai matriisien vaakasuuntainen ketjuttaminen} \\
  &|S| &&\text{joukon tai sekvenssin } S \text{ alkioiden lukumäärä} \\
  &\mathbf{1}\{\cdot\} &&\text{indikaattorifunktio} \\
  &\mathcal{V} = \{v_1, \dots, v_{|\mathcal{V}|}\} &&\text{sanasto (kaikkien käytettävien tokenien joukko)} \\
  &w_t \in \mathcal{V} &&\text{korpuksen ajanhetken } t \text{ tokeni} \\
  &\mathcal{W} = (w_1, \dots, w_{|\mathcal{W}|}) &&\text{korpus (havaittu tokenijono)} \\
  &x(v_j) \in \{0, 1\}^{1 \times |\mathcal{V}|} &&\text{tokenin } v_j \text{ one-hot-esitys (rivivektori)} \\
  &e(v_j) \in \mathbb{R}^{1 \times d} &&\text{tokenin } v_j \text{ tiheä upotusvektori (rivivektori)} \\
  &E \in \mathbb{R}^{|\mathcal{V}| \times d} &&\text{tokenien upotusmatriisi} \\
  &H \in \mathbb{R}^{T \times d} &&\text{esitysmatriisi (jonon pituus } T\text{, upotusdimensio } d\text{)} \\
  &\ell \in \mathbb{R}^{1 \times |\mathcal{V}|} &&\text{normalisoimaton pistemäärävektori (logittivektori)} \\
  &J(\theta) &&\text{häviöfunktio} \\
  &\theta &&\text{mallin opittavien parametrien joukko} \\
  &\Delta(\mathcal{V}) &&\text{sanaston } \mathcal{V} \text{ yli määriteltyjen todennäköisyysjakaumien joukko}
\end{align*}
